%-----------------------------------------------------------------------
% A248044_rev5.tex — Revision 5
%
% Computational Verification of Sun's Conjecture 4.3(i)
% and Exceptional Structure of Sequence A248044
%
% Rev. 5 changes (from 7-reviewer panel adjudication by Claude Opus 4.6):
%   [OBS-002] Shadowing discussion: added a(m)>1 guard (ChatGPT)
%   [OBS-003] CRITICAL: distribution table recomputed from CSV data
%   [OBS-004] Lemma 2.1 proof: more specific Ireland & Rosen citation (Kimi)
%   [OBS-006] Added Machiavelo & Tsopanidis (2021) reference (ChatGPT)
%   [OBS-007] Softened "first large-scale" claim in conclusions (ChatGPT)
%   [OBS-008] Added "heuristic" qualifiers to probability estimates (ChatGPT)
%   [OBS-011] More precise ratio in abstract (Grok)
%   [OBS-014] Added total CPU-hours (Gemini)
%   [OBS-015] Source of 57% elimination rate clarified (DeepSeek/Sonnet)
%
% Classe di lavoro: amsart (compila con qualsiasi TeX Live standard).
% Per la submission finale a MCOM, cambiare SOLO la prima riga:
%   \documentclass{amsart}  -->  \documentclass{mcom-l}
% e copiare mcom-l.cls nella cartella del progetto.
%-----------------------------------------------------------------------

\documentclass{amsart}

%--- Pacchetti ----------------------------------------------------------
\usepackage{amsmath,amssymb}
\usepackage{hyperref}
\usepackage{booktabs}
\usepackage{url}

%--- Configurazione hyperref --------------------------------------------
\hypersetup{
    colorlinks = true,
    linkcolor  = black,
    citecolor  = black,
    urlcolor   = black
}

%--- Ambienti teorema ---------------------------------------------------
\newtheorem{theorem}{Theorem}[section]
\newtheorem{lemma}[theorem]{Lemma}
\newtheorem{corollary}[theorem]{Corollary}
\newtheorem{proposition}[theorem]{Proposition}

\theoremstyle{definition}
\newtheorem{definition}[theorem]{Definition}
\newtheorem{example}[theorem]{Example}
\newtheorem{conjecture}[theorem]{Conjecture}
\newtheorem{openproblem}[theorem]{Open Problem}

\theoremstyle{remark}
\newtheorem{remark}[theorem]{Remark}

\numberwithin{equation}{section}

%--- Macro --------------------------------------------------------------
\newcommand{\N}{\mathbb{N}}
\newcommand{\Z}{\mathbb{Z}}
\newcommand{\Zp}{\mathbb{Z}^{+}}
\newcommand{\vq}[1]{v_{#1}}

%=======================================================================
\begin{document}
%=======================================================================

\title[Computational Verification of Sun's Conjecture~4.3(i)]
{Computational Verification of Sun's Conjecture~4.3(i)\\
and Exceptional Structure of Sequence~A248044}

\author[C. Corti]{Carlo Corti}
\address{Independent researcher, 21-040 Kalin\'owka, Poland}
\email{carlo.corti@outlook.com}

\thanks{%
  \textit{AI Disclosure.}
  The computational implementation, including the C~search programs
  and segmented sieve, was developed with extensive assistance
  from Anthropic's Claude language models
  (\texttt{claude-sonnet-4-6} and \texttt{claude-opus-4-6}).
  Claude also assisted substantially in the statistical analysis,
  the code review process, and in the preparation of this manuscript.
  All computations were executed and independently verified by the
  author, who bears full scientific responsibility for all results,
  including those portions where AI assistance was used.
  The use of AI tools is disclosed in accordance with AMS publication
  policy; see
  \url{https://www.ams.org/publications/journals/policies/UseofArtificialIntelligence}.%
}

\subjclass[2020]{Primary 11A41, 11Y55; Secondary 11N05, 11A07}

\keywords{prime counting function, sum of two squares,
computational verification, OEIS A248044, Sun's conjecture,
obstruction filter, exceptional cases}

\date{}

\begin{abstract}
We present the first large-scale computational investigation of
OEIS sequence A248044, which encodes Sun's Conjecture~4.3(i):
for every positive integer $m$, there exists a positive integer~$n$
such that $(m+n) \mid \pi(m)^2 + \pi(n)^2$, where $\pi(x)$ denotes
the prime-counting function.
We determine $a(m) = \min\{n \in \Zp : (m+n) \mid \pi(m)^2 + \pi(n)^2\}$
for $99\,998$ of the first $100\,000$ positive integers, with a
final search bound of $n \le 10^{13}$.
The computation employs a segmented sieve of Eratosthenes coupled
with a modular obstruction filter based on quadratic residuacity of~$-1$,
parallelised via OpenMP on~$16$ threads.
The two unresolved values, $m = 19\,623$ and $m = 19\,624$, satisfy
% REVISED [OBS-011]: more precise lower bound on ratio
$a(m) > 10^{13}$, yielding Cram\'er-type ratios
$a(m)/m > 5.096 \times 10^8$ --- the largest lower bounds in the
entire sequence.
We prove that the obstruction filter is sound
(Lemma~\ref{lem:obstruction}) and characterise the arithmetic
structure of the unresolved cases: all belong to a cluster sharing
$\pi(m) = 2225 = 5^2 \times 89$, whose prime factorisation consists
entirely of primes congruent to $1 \pmod{4}$, maximising the
modular obstruction.
Statistical analysis reveals a heavy-tailed distribution
consistent with a log-exponential model, and we provide
probabilistic estimates for the unresolved values.
\end{abstract}

\maketitle

%=======================================================================
\section{Introduction}
\label{sec:intro}
%=======================================================================

In a 2013 preprint subsequently published in 2019,
Zhi-Wei Sun~\cite{Sun2019} proposed a collection of conjectures
involving divisibility conditions on prime-related functions.
Among these, Conjecture~4.3(i) (p.~15 of~\cite{Sun2019}) states:

\begin{conjecture}[Sun, 2013/2019]\label{conj:sun43i}
For every positive integer $m$, there exists a positive integer $n$
such that
\begin{equation}\label{eq:divisibility}
(m + n) \mid \pi(m)^2 + \pi(n)^2,
\end{equation}
where $\pi(x) = \#\{p \le x : p \text{ prime}\}$ denotes the
prime-counting function.
\end{conjecture}

The sequence $a(m)$, defined as the least positive integer~$n$
satisfying~\eqref{eq:divisibility}, was submitted to the OEIS in
2014 and appears as sequence A248044~\cite{OEIS_A248044}.
The b-file at the time of writing contains values only for
$m \le 200$.
To the best of our knowledge, no large-scale computational
verification of Conjecture~\ref{conj:sun43i} has appeared in the
literature (for a related computational study of a different Sun
conjecture, see~\cite{MacTsop2021}).

The companion conjecture, 4.1(i), involving $p_m$ (the $m$-th prime)
in place of $\pi(m)$, was investigated computationally
in~\cite{CortiA247975}; the present paper extends the same
methodology to the prime-counting variant.

The present work makes three main contributions.

\emph{First}, we extend the computation of $a(m)$ to
$m = 1, \ldots, 100\,000$, resolving $99\,998$ of $100\,000$ cases
with a final search bound of $n \le 10^{13}$.  This is the largest
verified dataset for this conjecture.

\emph{Second}, we carry out a systematic statistical analysis of the
distribution of $a(m)$, revealing a heavy-tailed structure well
described by a log-exponential model.

\emph{Third}, we identify and characterise the two values of~$m$ for
which the conjecture resists verification beyond $n = 10^{13}$,
both belonging to an arithmetically structured cluster sharing a
common value $\pi(m) = 2225 = 5^2 \times 89$; we explain how the
factorisation of $\pi(m)$ into primes $\equiv 1 \pmod{4}$ produces
maximal modular obstruction and propose the anomalous resistance of
these two cases as an open problem.

The paper is organised as follows.
Section~\ref{sec:background} reformulates the divisibility condition
in modular arithmetic and establishes the theoretical basis for the
obstruction filter.
Section~\ref{sec:algorithm} describes the algorithm and implementation.
Section~\ref{sec:results} presents the computational results.
Section~\ref{sec:statistics} carries out the statistical analysis.
Section~\ref{sec:exceptional} analyses the two unresolved cases.
Section~\ref{sec:conclusions} concludes with open problems.

%=======================================================================
\section{Mathematical Background}
\label{sec:background}
%=======================================================================

For positive integers $m$ and $n$, set $s = m + n$.
The divisibility condition~\eqref{eq:divisibility} is equivalent to
\begin{equation}\label{eq:modular}
\pi(n)^2 \equiv -\pi(m)^2 \pmod{s}.
\end{equation}
When $\gcd(\pi(m), s) = 1$, this simplifies to
\begin{equation}\label{eq:qr}
\bigl(\pi(n) \cdot \pi(m)^{-1}\bigr)^2 \equiv -1 \pmod{s},
\end{equation}
so a solution exists at index~$n$ only if $-1$ is a quadratic
residue modulo~$s$ (a necessary but not sufficient condition, since
$\pi(n)$ must also fall in the correct residue class).

A classical result (see \cite{IrelandRosen}, Ch.~5) gives:

\begin{proposition}\label{prop:qr}
$-1$ is a quadratic residue modulo $s > 1$ if and only if both
of the following hold:
\begin{enumerate}
\item[\textup{(a)}] $4 \nmid s$
\textup{(}equivalently, $\vq{2}(s) \le 1$,
where $\vq{2}(s)$ denotes the $2$-adic valuation of~$s$\textup{)};
\item[\textup{(b)}] no prime $q \equiv 3 \pmod{4}$ divides~$s$.
\end{enumerate}
\end{proposition}

\begin{remark}\label{rem:jacobi}
Condition~(b) is strictly stronger than requiring the Jacobi
symbol $\bigl(\frac{-1}{s}\bigr) = 1$.
For odd~$s$, the Jacobi symbol equals~$1$ if and only if
$\sum_{q \equiv 3\,(4)} \vq{q}(s)$ is even; actual solvability of
$x^2 \equiv -1 \pmod{s}$ requires that no such prime divides~$s$
at all.
A standard example: $s = 21 = 3 \times 7$, where
$\bigl(\frac{-1}{21}\bigr) = 1$ yet $-1$ is not a QR modulo~$21$.
\end{remark}

\subsection{The obstruction filter}

We now formalise the obstruction that arises when a small prime
$q \equiv 3 \pmod{4}$ divides $s = m + n$ to an odd power but
does not divide~$\pi(m)$.

\begin{lemma}[Obstruction Filter]\label{lem:obstruction}
Let $q$ be a prime with $q \equiv 3 \pmod{4}$, and let
$s = m + n$ with $m, n \in \Zp$.
Suppose that $q \mid s$, that $\vq{q}(s)$ is odd, and that
$q \nmid \pi(m)$.
Then $s \nmid \pi(m)^2 + \pi(n)^2$.
\end{lemma}

\begin{proof}
Suppose for contradiction that $s \mid \pi(m)^2 + \pi(n)^2$.
Since $q \mid s$ and $\vq{q}(s)$ is odd, set $e = \vq{q}(s) \ge 1$ (odd).
Then $q^e \mid \pi(m)^2 + \pi(n)^2$.

Since $q \nmid \pi(m)$ by hypothesis, we have
$\gcd(\pi(m), q^e) = 1$, so $\pi(m)$ is invertible modulo~$q^e$.
From $q^e \mid \pi(m)^2 + \pi(n)^2$ we get
\[
\bigl(\pi(n) \cdot \pi(m)^{-1}\bigr)^2 \equiv -1 \pmod{q^e}.
\]
Reducing modulo~$q$:
\[
\bigl(\pi(n) \cdot \pi(m)^{-1}\bigr)^2 \equiv -1 \pmod{q}.
\]
Since $q \equiv 3 \pmod{4}$, the element $-1$ is not a quadratic
residue modulo~$q$ (\cite{IrelandRosen}, Ch.~5, Theorem~3).
This contradicts the equation above, so the assumption
$s \mid \pi(m)^2 + \pi(n)^2$ must be false.
\end{proof}

\begin{corollary}\label{cor:maxobstruction}
If every prime factor of $\pi(m)$ is congruent to $1 \pmod{4}$
\textup{(}i.e., $\pi(m)$ has no prime factor $q \equiv 3 \pmod{4}$\textup{)},
then for any small prime $q \equiv 3 \pmod{4}$ with $q \mid s$
and $\vq{q}(s)$ odd, the condition of
Lemma~\textup{\ref{lem:obstruction}} is satisfied and~$s$ is
blocked.
In this case the obstruction filter achieves its maximal
elimination rate.
\end{corollary}

\begin{proof}
If every prime factor of $\pi(m)$ is $\equiv 1 \pmod{4}$, then
$q \nmid \pi(m)$ for every $q \equiv 3 \pmod{4}$, and
Lemma~\ref{lem:obstruction} applies whenever $q \mid s$ with
$\vq{q}(s)$ odd.
\end{proof}

\begin{remark}\label{rem:filter_general}
The proof of Lemma~\ref{lem:obstruction} is valid for any prime
$q \equiv 3 \pmod{4}$, not only for the small primes used in our
implementation.  Our code tests $q \in \{3, 7, 11, 19\}$ for
computational efficiency; using all primes $q \equiv 3 \pmod{4}$
up to~$s$ would make the filter exact for condition~(b) of
Proposition~\ref{prop:qr} but at prohibitive computational cost.
\end{remark}

\subsection{Density of obstructed candidates}

The fraction of integers $s \le x$ that are divisible by at
least one prime $q \equiv 3 \pmod{4}$ to an odd power can be
estimated using the Landau--Ramanujan theorem~\cite{Landau1908}
and Mertens-type products~\cite{Mertens1874,Tenenbaum1995}.
The set of integers with no prime factor $q \equiv 3 \pmod{4}$
to odd power has natural density zero, with asymptotic count
$\sim K x / \sqrt{\ln x}$ for a constant $K \approx 0.7642$.
At the finite scales relevant to our computation ($x \sim 10^{9}$
to $10^{13}$), the fraction of candidate values of~$s$ that
\emph{pass} the full condition~(b) is roughly $17$--$20\%$;
the complementary $80$--$83\%$ are obstructed in principle.

Our implementation checks only $q \in \{3, 7, 11, 19\}$,
% REVISED [OBS-015]: clarify source of empirical elimination rate
which empirically eliminates approximately $57\%$ of candidates
(measured across the monolithic run for $m = 1, \ldots, 100\,000$),
well below the theoretical maximum.  The gap is due to the
limited trial-division bound: candidates whose only factor
$q \equiv 3 \pmod{4}$ exceeds~$19$ are not eliminated.

%=======================================================================
\section{Algorithm and Implementation}
\label{sec:algorithm}
%=======================================================================

\subsection{Overview}

The computation proceeded in two phases:
a \emph{monolithic} phase computing $a(m)$ for all
$m = 1, \ldots, 100\,000$ with search bound $n \le 6 \times 10^9$,
followed by a series of \emph{targeted segmented} searches
extending the bound to $n = 10^{13}$ for the unresolved cases.

All computations were performed on a personal workstation:
AMD Ryzen~9 7940HS (16 threads), 64\,GB DDR5 RAM,
running Linux Mint~22.3.  Compilation used GCC with
\texttt{-O3 -march=znver4 -fopenmp}.

\subsection{Monolithic search: \texttt{sun\_A248044\_v5.c}}

For the initial phase, the program builds a global table of
$\pi(n)$ for all $n \le N_{\max}$, stored as \texttt{uint32\_t}.
The sieve of Eratosthenes marks composites in a
\texttt{uint8\_t} array; a prefix sum then fills the $\pi$-table.
Memory consumption is $5 \cdot N_{\max}$ bytes
($\approx 30$\,GB for $N_{\max} = 6 \times 10^9$).

For each $m$ and each candidate $n$, the program sets
$s = m + n$ and checks:
\begin{enumerate}
\item \textbf{Obstruction filter:}
  for each $q \in \{3, 7, 11, 19\}$, if $q \mid s$, $\vq{q}(s)$
  is odd, and $q \nmid \pi(m)$, reject~$n$.
  Also reject if $4 \mid s$ (condition~(a) of
  Proposition~\ref{prop:qr}).
\item \textbf{Divisibility test:}
  compute $\pi(m)^2 + \pi(n)^2$ and check whether $s$ divides the
  sum.
\end{enumerate}

For $n \le 6 \times 10^9$, we have $\pi(n) \le \pi(6 \times 10^9)
\approx 2.80 \times 10^8$ and $\pi(m) \le \pi(100\,000)
= 9\,592$, so $\pi(m)^2 + \pi(n)^2 < 7.8 \times 10^{16}$, which
fits in \texttt{int64\_t} (max $\approx 9.2 \times 10^{18}$).

The monolithic phase resolved $99\,935$ of $100\,000$ cases in
approximately $10$~minutes of wall-clock time, leaving $65$
unresolved cases (the ``exceptional'' values of~$m$).

\subsection{Targeted segmented search: \texttt{sun\_A248044\_targeted\_v3.c}}
\label{sec:targeted}

The monolithic architecture cannot scale beyond
$\sim 7 \times 10^9$ due to the RAM ceiling of the global
$\pi$-table.
For the extended searches we developed a segmented approach:
the range $[n_{\min}, n_{\max})$ is divided into segments of width
$W = 10^9$, and for each segment $[l, l+W)$:

\begin{enumerate}
\item A segmented sieve of Eratosthenes marks composites in
  a \texttt{uint8\_t} array \texttt{seg[]} of size~$W$
  ($\approx 1$\,GB).
\item A prefix sum computes $\pi_{\mathrm{seg}}[k]
  = \#\{\text{primes in } [l, l+k]\}$ as \texttt{uint32\_t}
  ($\approx 4$\,GB), and the absolute count is recovered as
  $\pi(l+k) = \pi_{\mathrm{base}} + \pi_{\mathrm{seg}}[k]$,
  where $\pi_{\mathrm{base}} = \pi(l)$ is maintained across
  segments.
\item The inner loop iterates over all $k \in [0, W)$ in parallel
  (OpenMP, \texttt{schedule(guided)}), and for each still-unresolved
  target~$m$, applies the obstruction filter and divisibility test.
\end{enumerate}

Total memory per segment: $\approx 5$\,GB.  Since $\pi_{\mathrm{seg}}[k]$
counts primes within a single segment of width $10^9$, the offset
cannot exceed $\sim 4 \times 10^7$ and comfortably fits in
\texttt{uint32\_t}.

For $n > \sim 3.5 \times 10^9$ we have $\pi(n) > 1.6 \times 10^8$,
so $\pi(n)^2 > 2.6 \times 10^{16}$; and at the extremes of Ext-5
($n \sim 10^{13}$) we have $\pi(n) \approx 3.5 \times 10^{11}$, giving
$\pi(n)^2 \approx 1.2 \times 10^{23}$, which far exceeds the range
of \texttt{int64\_t}.  Accordingly, the segmented tool uses
\texttt{\_\_int128} (a GCC extension) uniformly for the sum
$\pi(m)^2 + \pi(n)^2$ and the modular reduction.
On x86-64, GCC compiles $\mathtt{\_\_int128} \mathbin{\%}
\mathtt{int64\_t}$ to a native \texttt{DIV} instruction
(\texttt{rdx:rax}), so the overhead relative to 64-bit arithmetic
is negligible (a few cycles per division).

\subsection{Parallelism and correctness}

The inner loop is parallelised with
\texttt{\#pragma omp parallel for schedule(guided)}.
The choice of \texttt{guided} scheduling reflects the non-uniform
cost per iteration: the obstruction filter rejects $\sim 57\%$ of
candidates without performing the expensive 128-bit division,
creating significant load imbalance under static scheduling.

Target data (arrays of $m$, $\pi(m)^2$, and filter bitmasks)
are copied into Structure-of-Arrays (SoA) layout before the
parallel region.  With at most $65$ targets at $\le 24$ bytes each,
the working set of $\approx 1.5$\,kB fits entirely in the L1 cache
(32\,kB per core), ensuring that the hot loop is memory-bound on
the segment arrays rather than on target metadata.

The variable \texttt{found\_n[ai]} (recording the best solution
found so far for target~$ai$) is read without synchronisation
inside the parallel loop and written only under
\texttt{\#pragma omp critical}.  This constitutes a formal data
race under the C11/OpenMP memory model; however, on x86-64 under
the Total Store Order (TSO) memory model, an aligned 64-bit read
cannot produce a torn value.  The worst case is a stale read, which
may cause one redundant iteration but never affects correctness:
the critical section always compares and retains the global minimum.
An \texttt{omp atomic read} directive would eliminate the formal
race but at a cost of $\sim 4$ cycles $\times 6.5 \times 10^{10}$
iterations $\approx 87$~seconds per segment --- unacceptable given
that the benign race preserves correctness on the target
architecture.

\subsection{Code review}

The targeted search code (v3) underwent two rounds of independent
review with assistance from AI language models (see the AI~Disclosure
footnote on p.~1).
The first round identified two critical bugs
(integer overflow in $\pi(n)^2$ and a race condition on the
minimum-tracking logic); both were fixed in v2.
The second round confirmed correctness and led to
minor improvements (input parsing hardening, overflow fix in the
CSV output column).
All observations and their dispositions are documented in the source
code header.

%=======================================================================
\section{Computational Results}
\label{sec:results}
%=======================================================================

\subsection{Campaign summary}

Table~\ref{tab:campaigns} summarises the progressive resolution of
the $65$ initially unresolved cases across the monolithic and
extended campaigns.

\begin{table}[ht]
\centering
\caption{Campaign summary.  The monolithic run (M1) used
\texttt{sun\_A248044\_v5.c}; all Ext runs used
\texttt{sun\_A248044\_targeted\_v3.c}.
``+Resolved'' denotes newly resolved cases in each run.}
\label{tab:campaigns}
\small
\begin{tabular}{llrrrl}
\toprule
Run & Bound $n$ & +Resolved & Still open & Time & Note \\
\midrule
M1 (monolithic) & $\le 6 \times 10^9$
  & $99\,935$ & 65 & $\sim$10\,min & All $m \le 10^5$ \\
Ext-2 & $6 \times 10^9 \to 10^{10}$
  & 36 & 29 & 137\,s & \\
Ext-3 & $10^{10} \to 10^{11}$
  & 20 & 9 & 1\,132\,s & \\
Ext-4 & $10^{11} \to 10^{12}$
  & 4 & 5 & 9\,452\,s & \\
Ext-5 & $10^{12} \to 10^{13}$
  & 3 & \textbf{2} & 86\,098\,s & \\
\bottomrule
\end{tabular}
\end{table}

The total computation time for the extended campaigns was
approximately $96\,819$ seconds ($\approx 26.9$ hours).
% REVISED [OBS-014]: explicit total wall-clock and CPU-hours
Including the monolithic phase ($\approx 0.17$~hours), the entire
computation required approximately $27$~hours of wall-clock time
($\approx 432$ CPU-hours on $16$ threads).
An Ext-6 campaign ($10^{13} \to 10^{14}$) would require an
estimated $\sim 240$ hours of wall-clock time on the same hardware
and was deemed infeasible.

\subsection{Largest values of $a(m)$}

Table~\ref{tab:top10_absolute} lists the ten largest absolute
values of $a(m)$ found in the computation.

\begin{table}[ht]
\centering
\caption{Top~10 values of $a(m)$ by absolute magnitude.}
\label{tab:top10_absolute}
\small
\begin{tabular}{rrrrl}
\toprule
$m$ & $\pi(m)$ & $a(m)$ & Ratio $a(m)/m$ & Campaign \\
\midrule
$83\,115$ & $8\,117$ & $4\,629\,736\,117\,663$ & $5.57 \times 10^7$ & Ext-5 \\
$40\,963$ & $4\,289$ & $2\,650\,225\,478\,854$ & $6.47 \times 10^7$ & Ext-5 \\
$19\,620$ & $2\,225$ & $1\,150\,315\,774\,321$ & $5.86 \times 10^7$ & Ext-5 \\
$57\,919$ & $5\,867$ & $330\,283\,186\,011$ & $5.70 \times 10^6$ & Ext-4 \\
$57\,920$ & $5\,867$ & $330\,283\,186\,010$ & $5.70 \times 10^6$ & Ext-4 \\
$45\,151$ & $4\,687$ & $229\,736\,032\,374$ & $5.09 \times 10^6$ & Ext-4 \\
$93\,381$ & $9\,019$ & $190\,261\,921\,669$ & $2.04 \times 10^6$ & Ext-4 \\
$85\,504$ & $8\,321$ & $79\,705\,931\,373$ & $9.32 \times 10^5$ & Ext-3 \\
$44\,857$ & $4\,661$ & $32\,348\,326\,705$ & $7.21 \times 10^5$ & Ext-3 \\
$44\,858$ & $4\,661$ & $32\,348\,326\,704$ & $7.21 \times 10^5$ & Ext-3 \\
\bottomrule
\end{tabular}
\end{table}

\subsection{Largest Cram\'er-type ratios}

Table~\ref{tab:top10_ratio} lists the ten largest
Cram\'er-type ratios $C(m) = a(m)/m$.

\begin{table}[ht]
\centering
\caption{Top~10 values by Cram\'er-type ratio $a(m)/m$.}
\label{tab:top10_ratio}
\small
\begin{tabular}{rrrrr}
\toprule
$m$ & $\pi(m)$ & $a(m)$ & $a(m)/m$ & Campaign \\
\midrule
$40\,963$ & $4\,289$ & $2\,650\,225\,478\,854$ & $6.47 \times 10^7$ & Ext-5 \\
$19\,620$ & $2\,225$ & $1\,150\,315\,774\,321$ & $5.86 \times 10^7$ & Ext-5 \\
$83\,115$ & $8\,117$ & $4\,629\,736\,117\,663$ & $5.57 \times 10^7$ & Ext-5 \\
$57\,919$ & $5\,867$ & $330\,283\,186\,011$ & $5.70 \times 10^6$ & Ext-4 \\
$57\,920$ & $5\,867$ & $330\,283\,186\,010$ & $5.70 \times 10^6$ & Ext-4 \\
$45\,151$ & $4\,687$ & $229\,736\,032\,374$ & $5.09 \times 10^6$ & Ext-4 \\
$93\,381$ & $9\,019$ & $190\,261\,921\,669$ & $2.04 \times 10^6$ & Ext-4 \\
$85\,504$ & $8\,321$ & $79\,705\,931\,373$ & $9.32 \times 10^5$ & Ext-3 \\
$44\,857$ & $4\,661$ & $32\,348\,326\,705$ & $7.21 \times 10^5$ & Ext-3 \\
$44\,858$ & $4\,661$ & $32\,348\,326\,704$ & $7.21 \times 10^5$ & Ext-3 \\
\bottomrule
\end{tabular}
\end{table}

\subsection{Distribution by order of magnitude}

Table~\ref{tab:distribution} shows the distribution of $a(m)$
across orders of magnitude, among the $99\,998$ resolved cases.

% REVISED [OBS-003]: CRITICAL FIX — distribution table recomputed from CSV data.
% Original table had fabricated counts; these are verified from
% out_A248044_m100000.csv + ext CSVs + 149 values from M1 6e9 run.
\begin{table}[ht]
\centering
\caption{Distribution of $a(m)$ by order of magnitude
($99\,998$ resolved cases plus $2$ unresolved).}
\label{tab:distribution}
\begin{tabular}{lrr}
\toprule
Range & Count & Cumulative \\
\midrule
$a(m) \le 10^6$     & $89\,685$ & $89\,685$ \\
$10^6 < a(m) \le 10^7$   & $7\,516$ & $97\,201$ \\
$10^7 < a(m) \le 10^8$   & $2\,009$ & $99\,210$ \\
$10^8 < a(m) \le 10^9$   & $576$ & $99\,786$ \\
$10^9 < a(m) \le 10^{10}$  & $185$ & $99\,971$ \\
$10^{10} < a(m) \le 10^{11}$ & $20$ & $99\,991$ \\
$10^{11} < a(m) \le 10^{12}$ & $4$ & $99\,995$ \\
$10^{12} < a(m) \le 10^{13}$ & $3$ & $99\,998$ \\
$a(m) > 10^{13}$ (unresolved) & $2$ & $100\,000$ \\
\bottomrule
\end{tabular}
\end{table}

The monotonically decreasing counts across orders of magnitude
confirm the heavy-tailed nature of the distribution.
% REVISED [OBS-003]: corrected median from actual CSV data
The median value is $a(m) = 52\,846$.

\subsection{Independent verification}

All $63$ results from the extended campaigns (Ext-2 through Ext-5)
were independently verified by checking the divisibility condition
\begin{equation}\label{eq:verify}
(\pi(m)^2 + \pi(a(m))^2) \bmod (m + a(m)) = 0
\end{equation}
using the values $\pi(m)$ and $\pi(a(m))$ reported in the CSV
output files.  These values can be verified against the known
checkpoints $\pi(10^k)$ from OEIS A006880~\cite{OEIS_A006880};
Table~\ref{tab:pi_checkpoints} lists those used.

\begin{table}[ht]
\centering
\caption{Checkpoints for $\pi(n)$ used in independent verification
(OEIS A006880, confirmed with SymPy~\cite{SymPy}).}
\label{tab:pi_checkpoints}
\begin{tabular}{rr}
\toprule
$n$ & $\pi(n)$ \\
\midrule
$10^{8}$ & $5\,761\,455$ \\
$10^{9}$ & $50\,847\,534$ \\
$6 \times 10^{9}$ & $279\,545\,368$ \\
$10^{10}$ & $455\,052\,511$ \\
$10^{11}$ & $4\,118\,054\,813$ \\
$10^{12}$ & $37\,607\,912\,018$ \\
$10^{13}$ & $346\,065\,536\,839$ \\
\bottomrule
\end{tabular}
\end{table}

\subsection{Consecutive-$m$ pairs}

% REVISED [OBS-002]: added a(m)>1 guard and edge-case note
A striking feature of the dataset is the prevalence of consecutive
pairs $(m, m+1)$ sharing large values of $a$.  The top-10 tables
include the pairs $(57\,919, 57\,920)$ and $(44\,857, 44\,858)$,
each with $a(m+1) = a(m) - 1$.  This is no coincidence: provided
$a(m) > 1$, if both $\pi(m+1) = \pi(m)$ (i.e., $m+1$ is composite)
and $\pi(a(m)-1) = \pi(a(m))$ (i.e., $a(m)$ is composite), then
$s' = (m+1) + (a(m)-1) = m + a(m) = s$, and the divisibility
$s \mid \pi(m)^2 + \pi(a(m))^2$ immediately yields
$s' \mid \pi(m+1)^2 + \pi(a(m)-1)^2$, so
$a(m+1) \le a(m) - 1$.
(When $a(m) = 1$, as happens for $m = 3, 7, \ldots$, the candidate
$a(m) - 1 = 0$ is not a positive integer and the mechanism does
not apply.)
This ``shadowing'' mechanism is not guaranteed (either $m+1$ or
$a(m)$ could be prime), but when it holds, it produces
exact pairs with $a(m+1) = a(m) - 1$.

%=======================================================================
\section{Statistical Analysis}
\label{sec:statistics}
%=======================================================================

\subsection{Heavy-tailed distribution}

The distribution of $a(m)$ is strongly heavy-tailed: while the
% REVISED [OBS-003]: corrected median
median is approximately $52\,846$, the maximum resolved value is
$\approx 4.63 \times 10^{12}$, a ratio of $\sim 8.8 \times 10^7$.
Values of $a(m)$ exceeding $10^9$ occur for $63$ values of~$m$
out of $99\,998$ resolved ($0.063\%$), while values exceeding
$10^{12}$ occur for $3$ values ($0.003\%$).

\subsection{A log-exponential model}

We model the occurrence of solutions heuristically.
For a given~$m$, each candidate $s = m + n$ (not blocked by the
obstruction filter) independently has a ``pass probability'' that
depends on the number of valid residue classes for $\pi(n)$
modulo~$s$.  Under a Cram\'er-type heuristic, the number of
solutions in an interval $[N_1, N_2]$ is approximately proportional
to $\ln(N_2/N_1)$, giving a log-exponential distribution:
\begin{equation}\label{eq:logexp}
\Pr\bigl(a(m) < N \mid a(m) > N_0\bigr)
= 1 - \exp\bigl(-\lambda \cdot \ln(N/N_0)\bigr)
= 1 - (N_0/N)^{\lambda},
\end{equation}
where $\lambda > 0$ is a pass-rate parameter.

\subsection{Calibration from extended campaigns}

The pass-rate parameter~$\lambda$ can be estimated from the
extended campaigns by the relation
\[
\lambda \approx
\frac{\text{(number resolved)}}
     {\text{(number of targets)} \times \ln(N_{\max}/N_{\min})}.
\]
Table~\ref{tab:passrate} summarises the estimates.

\begin{table}[ht]
\centering
\caption{Empirical pass rates by campaign.  The declining
trend reflects the selection effect: cases surviving to higher
bounds have intrinsically lower pass rates.}
\label{tab:passrate}
\begin{tabular}{lrrl}
\toprule
Campaign & Resolved/Targets & $\ln(N_{\max}/N_{\min})$ &
  $\hat\lambda$ \\
\midrule
Ext-2 & $36/65$  & $0.511$ & $1.08$ \\
Ext-3 & $20/29$  & $2.303$ & $0.30$ \\
Ext-4 & $4/9$    & $2.303$ & $0.19$ \\
Ext-5 & $3/5$    & $2.303$ & $0.26$ \\
\bottomrule
\end{tabular}
\end{table}

The declining pass rate from Ext-2 to the later campaigns
reflects a strong selection effect: the cases surviving to higher
bounds are precisely those with the highest intrinsic obstruction,
so the average pass rate among survivors decreases with the bound.

Using the Ext-5 estimate $\hat\lambda \approx 0.26$ for the
two unresolved cases, the conditional median (given $a(m) > 10^{13}$)
is
\[
\tilde{N} = 10^{13} \cdot \exp(\ln 2 / 0.26)
\approx 10^{13} \times 14.4
\approx 1.4 \times 10^{14}.
\]

\subsection{Probabilistic estimates for unresolved cases}

% REVISED [OBS-008]: explicit heuristic qualifiers on probability statements
Under the log-exponential model with $\lambda = 0.26$,
the \emph{heuristic} probability of resolving each unresolved case
within various bounds is:
\begin{align*}
\Pr\bigl(a(m) \le 10^{14} \mid a(m) > 10^{13}\bigr)
  &\approx 1 - 10^{-0.26} \approx 0.45, \\
\Pr\bigl(a(m) \le 10^{15} \mid a(m) > 10^{13}\bigr)
  &\approx 1 - 10^{-0.52} \approx 0.70, \\
\Pr\bigl(a(m) \le 10^{16} \mid a(m) > 10^{13}\bigr)
  &\approx 1 - 10^{-0.78} \approx 0.83.
\end{align*}

The heuristic probability of finding \emph{both} cases within
$[10^{13}, 10^{14}]$ is approximately $0.45^2 \approx 0.20$,
and within $[10^{13}, 10^{15}]$ approximately
$0.70^2 \approx 0.49$.

\begin{remark}
The log-exponential model treats $a(m)$ as though each logarithmic
unit of search space is equally likely to yield a solution.  This is
a heuristic, not a theorem; the data are deterministic, not i.i.d.,
and the model does not account for the specific arithmetic structure
of individual cases.  The estimates above should be understood as
order-of-magnitude guidance, not precise probabilities.
\end{remark}

%=======================================================================
\section{Exceptionally Resistant Cases}
\label{sec:exceptional}
%=======================================================================

\subsection{The two unresolved values}

After the full campaign (bound $n = 10^{13}$), the two values
\begin{equation}\label{eq:open}
m = 19\,623 \quad\text{and}\quad m = 19\,624
\end{equation}
remain unresolved, with
\[
a(19\,623) > 10^{13}, \qquad a(19\,624) > 10^{13},
\]
giving Cram\'er-type lower bounds
% REVISED [OBS-011]: precise lower bound consistent with abstract
$a(m)/m > 5.096 \times 10^8$ --- the largest in the entire dataset.

\subsection{The cluster $\{19\,620, 19\,623, 19\,624\}$}

The three values $m \in \{19\,620, 19\,623, 19\,624\}$ all satisfy
$\pi(m) = 2225 = 5^2 \times 89$.  The value $m = 19\,620$ was
resolved in Ext-5 at $a(19\,620) = 1\,150\,315\,774\,321$
(ratio $\approx 5.86 \times 10^7$), but $m = 19\,623$ and
$m = 19\,624$ resist beyond $10^{13}$ --- a factor of at least
$8.7$ further.

The factorisation $2225 = 5^2 \times 89$ consists entirely of
primes congruent to $1 \pmod{4}$ ($5 \equiv 1 \pmod{4}$ and
$89 \equiv 1 \pmod{4}$).  By Corollary~\ref{cor:maxobstruction},
the obstruction filter achieves its maximal elimination rate for
all three values, since no filter prime $q \in \{3, 7, 11, 19\}$
(all $\equiv 3 \pmod{4}$) divides $\pi(m) = 2225$.

This maximal obstruction explains why all three values are
exceptional, but it does not explain why $m = 19\,623$ and
$m = 19\,624$ resist so much more than $m = 19\,620$.

\subsection{Congruence analysis}

Table~\ref{tab:congruences} compares the residues of the three
cluster members modulo small numbers.

\begin{table}[ht]
\centering
\caption{Residues of the cluster $\{19\,620, 19\,623, 19\,624\}$
modulo small values.}
\label{tab:congruences}
\begin{tabular}{rrrrrrl}
\toprule
$m$ & $m \bmod 3$ & $m \bmod 4$ & $m \bmod 7$ & $m \bmod 11$
  & $m \bmod 19$ & Status \\
\midrule
$19\,620$ & 0 & 0 & 6 & 7 & 12 & Resolved ($1.15 \times 10^{12}$) \\
$19\,623$ & 0 & \textbf{3} & 2 & 10 & 15 & Open ($> 10^{13}$) \\
$19\,624$ & 1 & 0 & 3 & \textbf{0} & 16 & Open ($> 10^{13}$) \\
\bottomrule
\end{tabular}
\end{table}

Notable features:
\begin{itemize}
\item $m = 19\,623 \equiv 3 \pmod{4}$: for candidate
  $n \equiv 2 \pmod{4}$, $s = m + n \equiv 1 \pmod{4}$
  (favourable); but for $n \equiv 1 \pmod{4}$,
  $s \equiv 0 \pmod{4}$ (blocked by condition~(a)).
  This creates a more restrictive sieve pattern than for the other
  two values.
\item $11 \mid 19\,624$: since $11 \equiv 3 \pmod{4}$ and
  $11 \nmid \pi(19\,624) = 2225$, any $s$ with $\vq{11}(s)$ odd
  is blocked.  The constraint $11 \mid m$ forces
  $s = m + n \equiv n \pmod{11}$, creating a periodic pattern
  that may interact unfavourably with the other obstructions.
\end{itemize}

The precise mechanism by which these congruence patterns produce an
extra factor of $>8.7$ in resistance relative to $m = 19\,620$
remains unexplained, and we propose it as an open problem
(see Section~\ref{sec:conclusions}).

%=======================================================================
\section{Conclusions and Open Problems}
\label{sec:conclusions}
%=======================================================================

We have determined $a(m)$ for $99\,998$ of the first $100\,000$
positive integers, with final search bound $n \le 10^{13}$,
% REVISED [OBS-007]: softened "first" claim
providing, to the best of our knowledge, the first large-scale
computational verification of Sun's Conjecture~4.3(i).  The dataset reveals a heavy-tailed distribution
well described by a log-exponential model, with record values
reaching $a(83\,115) = 4\,629\,736\,117\,663$ and
Cram\'er-type ratios up to $6.47 \times 10^7$.

The progressive resolution of the exceptional cases --- from $65$
unresolved at bound $6 \times 10^9$, to $29$ at $10^{10}$, to~$9$
at $10^{11}$, to~$5$ at $10^{12}$, and finally to~$2$ at
$10^{13}$ --- strongly supports the conjecture: the pattern is
consistent with a heavy-tailed distribution where
every case is eventually resolved at a sufficiently high bound,
rather than with the existence of structural counterexamples.

We conclude with three open problems.

\begin{openproblem}\label{op:unresolved}
Determine $a(19\,623)$ and $a(19\,624)$.  Both satisfy
$a(m) > 10^{13}$, with $a(m)/m > 5.096 \times 10^8$.
The heuristic model suggests a median around
$1.4 \times 10^{14}$, with approximately~$70\%$ probability
of resolution within $n \le 10^{15}$.
An Ext-6 campaign ($10^{13} \to 10^{14}$) would require
$\sim 240$~hours on our hardware; algorithmic improvements
(e.g., bitset sieve with $4\times$ segment width,
extended filter with $q < 200$) might reduce this to
$\sim 35$--$40$~hours.
\end{openproblem}

\begin{openproblem}\label{op:cluster}
Explain the anomalous resistance of $m = 19\,623$ and
$m = 19\,624$ relative to $m = 19\,620$, all sharing
$\pi(m) = 2225$.  The resolved case $a(19\,620) \approx
1.15 \times 10^{12}$ is already extreme, yet the unresolved
cases exceed it by a factor of at least~$8.7$.
Is there an arithmetic characterisation of the congruence
conditions on~$m$ (beyond the factorisation of~$\pi(m)$) that
predicts this extra resistance?
\end{openproblem}

\begin{openproblem}\label{op:model}
Develop a refined probabilistic model for $a(m)$ that
accounts for the dependence on the factorisation of~$\pi(m)$.
The current log-exponential model treats all unresolved cases
uniformly; a model incorporating the obstruction rate
(which depends on the prime factorisation of~$\pi(m)$ and the
congruence class of~$m$) would provide sharper predictions.
\end{openproblem}

%=======================================================================
\section*{Data Availability}
%=======================================================================

The complete dataset ($99\,998$ resolved values and the
$2$~unresolved $m$-values), the C~source code for the monolithic
search (\texttt{sun\_A248044\_v5.c}) and the segmented targeted
search (\texttt{sun\_A248044\_targeted\_v3.c}), the full CSV
outputs for all campaigns, the machine-readable verification
logs, and a Python verification script (\texttt{verify\_A248044.py})
that checks the divisibility condition
$(m+a(m)) \mid \pi(m)^2 + \pi(a(m))^2$ for any supplied list
of $(m, a(m))$ pairs are publicly available at:
\begin{itemize}
\item GitHub repository:
      \url{https://github.com/carcorti/A248044}
\item Permanent Zenodo archive (DOI): to be assigned upon
      submission.
\end{itemize}
An extended b-file for OEIS A248044 has been submitted.

%=======================================================================
\section*{Acknowledgements}
%=======================================================================

The author thanks Zhi-Wei Sun for the beautiful conjecture that
motivated this work.
The computational implementation, including the C~search programs
and segmented sieve, was developed with extensive assistance from
Anthropic's Claude language models (\texttt{claude-sonnet-4-6} and
\texttt{claude-opus-4-6}); Claude also assisted substantially in
the statistical analysis, the code review process, and in the
preparation of the manuscript.
All computations were executed and independently verified by the
author, who bears full scientific responsibility for the results.

%=======================================================================
\bibliographystyle{amsplain}
\begin{thebibliography}{99}

\bibitem{Sun2019}
Zhi-Wei Sun,
\textit{Some new problems in additive combinatorics},
Nanjing Univ.\ J.\ Math.\ Biquarterly \textbf{36} (2019), no.~2,
134--155; preprint: arXiv:1309.1679 [math.NT] (submitted 2013,
v9: 2020).
Conjecture~4.3(i) appears on p.~15.
\url{https://arxiv.org/abs/1309.1679}

\bibitem{CortiA247975}
C.~Corti,
\textit{Exceptional cases and statistical structure of sequence
A247975: a computational study of Sun's Conjecture~4.1(i)},
preprint, 2026.

\bibitem{OEIS_A248044}
OEIS Foundation,
\textit{Sequence A248044},
The On-Line Encyclopedia of Integer Sequences
(accessed March~2026).
\url{https://oeis.org/A248044}

\bibitem{OEIS_A006880}
OEIS Foundation,
\textit{Sequence A006880} ($\pi(10^n)$),
The On-Line Encyclopedia of Integer Sequences
(accessed March~2026).
\url{https://oeis.org/A006880}

\bibitem{IrelandRosen}
K.~Ireland and M.~Rosen,
\textit{A Classical Introduction to Modern Number Theory},
2nd~ed., Springer, New York, 1990.

\bibitem{CrandallPomerance}
R.~Crandall and C.~Pomerance,
\textit{Prime Numbers: A Computational Perspective},
2nd~ed., Springer, New York, 2005.

\bibitem{Cramer1936}
H.~Cram\'er,
\textit{On the order of magnitude of the difference between
consecutive prime numbers},
Acta Arith.\ \textbf{2} (1936), 23--46.

\bibitem{Granville1995}
A.~Granville,
\textit{Harald Cram\'er and the distribution of prime numbers},
Scand.\ Actuarial J.\ (1995), no.~1, 12--28.

\bibitem{Landau1908}
E.~Landau,
\textit{\"Uber die Einteilung der positiven ganzen Zahlen in vier
Klassen nach der Mindestzahl der zu ihrer additiven Zusammensetzung
erforderlichen Quadrate},
Arch.\ Math.\ Phys.\ \textbf{13} (1908), 305--312.

\bibitem{Mertens1874}
F.~Mertens,
\textit{Ein Beitrag zur analytischen Zahlentheorie},
J.\ Reine Angew.\ Math.\ \textbf{78} (1874), 46--62.

\bibitem{Tenenbaum1995}
G.~Tenenbaum,
\textit{Introduction to Analytic and Probabilistic Number Theory},
Cambridge Univ.\ Press, 1995.

\bibitem{Hill1975}
B.~M.~Hill,
\textit{A simple general approach to inference about the tail of a
distribution},
Ann.\ Statist.\ \textbf{3} (1975), 1163--1174.

\bibitem{SymPy}
A.~Meurer et al.,
\textit{SymPy: symbolic computing in Python},
PeerJ Computer Science \textbf{3} (2017), e103.

\bibitem{OpenMP}
OpenMP Architecture Review Board,
\textit{OpenMP Application Programming Interface}, Version~5.2, 2021.
\url{https://www.openmp.org/specifications/}

% REVISED [OBS-006]: added contextual Sun-conjecture computational reference
\bibitem{MacTsop2021}
A.~Machiavelo and A.~Tsopanidis,
\textit{Zhi-Wei Sun's 1-3-5 conjecture and variations},
J.\ Number Theory \textbf{222} (2021), 135--155.

\end{thebibliography}

\end{document}
%-----------------------------------------------------------------------
% Per la submission finale a MCOM:
%   Cambiare \documentclass{amsart}
%         in \documentclass{mcom-l}
%   e copiare mcom-l.cls nella cartella del progetto.
% Tutto il resto rimane invariato.
%-----------------------------------------------------------------------
